\maketitle

\begin{abstract}
Recursive self-improvement claims are easy to overstate because improvements in a search process, a task solver, an evaluator, and a model's underlying capabilities are different phenomena. We report a bounded, prospectively validated two-transition process-improvement chain in \genesis, a persistent software-lineage experiment whose editable search policy is separated from a fixed external organism evaluator.

The chain begins with an offline meta-search result, V14, that produces a first child generation G1. In the subsequent real-organism process experiment, V15.4, G1 reaches the same retained champion quality as G0 while consuming two hidden-evaluator observations instead of four, yielding frozen lexicographic utilities G0=(0,0,-4) and G1=(0,0,-2). Attempts to continue recursively then fail in V16 and V17, and V18 exposes a specific overfitting shortcut: four independently proposed executable children all improve the public replay surface but tie G1 on all 24 fresh holdout worlds.

V19 responds prospectively by expanding development to a stratified public corpus of 13 real replay worlds and 188 synthetic worlds, while forbidding direct access to world identity. Four isolated proposer attempts converge to decision-equivalent policies. The selected executable G2 is frozen before a new run-ID-derived holdout exists. On 24 fresh worlds, G2 records 24 wins, 0 losses and 0 ties against G1, improving mean replay core by 75.25 while reducing represented requests by 1.0417. V20 then applies G1 and G2 to the same reconstructed real A016 organism state. G2 prospectively requests no hidden evaluation; G1 requests two physical evaluations, both of which are neutral at quality 953. Terminal utilities are therefore G1=(0,0,-2) and G2=(0,0,0).

Under the project's preregistered definition, these results establish a bounded \lthreep recursive process-improvement chain across G0$\rightarrow$G1$\rightarrow$G2. They do not establish recursive quality uplift at equal budget, transfer to a new evaluator or organism family, fully autonomous model-only self-authorship, or open-ended recursive self-improvement.
\end{abstract}

\section*{Claim in one sentence}
Under prospectively frozen rules, a lineage-mediated search policy improved twice across successive generations: first by reducing hidden evaluation cost at unchanged retained quality, then by producing a G2 policy that generalized on an unseen replay holdout and matched G1's real-organism outcome with zero hidden evaluations instead of two.

\section*{Evidence snapshot}
\begin{center}
\begin{tabular}{@{}llll@{}}
\toprule
Stage & Parent & Child & Frozen outcome \\
\midrule
V14 offline & G0 & G1 & 1 future win, 0 losses, 1 adoption \\
V15.4 real process & G0 & G1 & (0,0,-4) vs (0,0,-2) \\
V18 holdout & G1 & candidate G2 & 0 wins, 0 losses, 24 ties \\
V19 holdout & G1 & admitted G2 & 24 wins, 0 losses, 0 ties \\
V20 real process & G1 & G2 & (0,0,-2) vs (0,0,0) \\
\bottomrule
\end{tabular}
\end{center}
